% Partie : set-theory
% Chapitre : ordinals
% Section : milestone

\documentclass[../../../include/open-logic-section]{subfiles}

\begin{document}
\olfileid[fr]{sth}{ordinals}{zfm}
\olsection{\texorpdfstring{$\ZFminus$}{ZF⁻} : un jalon}

La question de savoir comment justifier le remplacement, si tant est
qu'on le puisse, n'est pas simple. Nous la réserverons donc au
\olref[replacement][]{chap}. Toutefois, avec l'ajout du
remplacement, nous avons franchi un autre jalon important. Nous
disposons désormais de tous les axiomes nécessaires à la théorie
$\ZFminus$. Plus précisément :

\begin{defn}
La théorie $\ZFminus$ a pour axiomes l'extensionnalité, l'union, la
paire, l'ensemble des parties, l'infini, ainsi que toutes les
instances des schémas de séparation et de remplacement. Autrement
dit, $\ZFminus$ ajoute le remplacement à $\Zminus$.
\end{defn}
\noindent
Cette abréviation désigne la théorie des ensembles de
\emph{Zermelo--Fraenkel} (\emph{moins} quelque chose dont nous
parlerons plus tard). Fraenkel a droit à cet honneur, car on lui
attribue la formulation du remplacement en
\citeyear{Fraenkel1922}, même si la première formulation précise est
due à \citet{Skolem1922}.

\end{document}
